% =============================================================================
%  CHAPTER 8 -- CONCLUSION                              [v3 drafts, v4 refines]
% -----------------------------------------------------------------------------
%  v3.7 (2026-09-14): shared between v3 and v4 (Working_Space/DRAFT_OWNERSHIP.md).
%  v3 keeps a concise version that follows Chapter 6; v4 refines the prose and
%  writes sec:conc:future. Individual changelog:
%  changelogs/v3.7_20260914_ch5-8_restructure.md
% =============================================================================

\chapter{Conclusion}\label{ch:conclusion}

\section{Summary}\label{sec:conc:summary}

This thesis replaced the diffusion model of \ac{DPCC} with instantaneous-velocity matching, analytic
average-velocity matching and consistency-interpolated average-velocity matching, keeping the backbone,
the projection, the data and the evaluation code of the baseline. The
new models were first evaluated on the D3IL-avoiding benchmark of \ac{DPCC} with its setup
unchanged, where the published baseline result was reproduced, and then on a vision-conditioned
D3IL-aligning task and a quadrotor benchmark built for this thesis.

On the benchmark of \ac{DPCC}, analytic average-velocity matching at a single network evaluation reaches the success with
constraint satisfaction of the baseline in fewer control steps and at a small fraction of its
wall-clock time, and its plans are coherent before projection where those of the diffusion model at the
same budget are not. On D3IL-aligning it brings the box closer to its target than the baseline and
instantaneous-velocity matching in every context, and on the quadrotor it completes UAV-corridor collision-free and leads on
UAV-pillars. Consistency-interpolated average-velocity matching is level with analytic average-velocity matching where the tasks cannot separate them and behind
it where they can. Endpoint projection is not used at the one- and two-evaluation operating point of
D3IL-avoiding. At larger budgets it reduces time on D3IL-aligning and UAV-corridor and projection time
on UAV-pillars, while its constraint result is higher on D3IL-aligning and UAV-pillars and lower on
UAV-corridor.

\hole{Refine in v4.}

\section{Answers to the Research Questions}\label{sec:conc:rqs}

\begin{description}
  \item[RQ1 --- Generative model.] With the backbone, projection, data and evaluation code held fixed,
    analytic average-velocity matching is ahead of instantaneous-velocity matching and of the diffusion baseline in every environment. Instantaneous-velocity matching
    is ahead of the baseline on D3IL-avoiding and UAV-pillars and level with it on D3IL-aligning. Consistency-interpolated average-velocity matching is not ahead of analytic average-velocity matching in
    any environment (\autoref{sec:res:summary:models}).
  \item[RQ2 --- Budget.] The flow-based models keep their success down to one network evaluation on the
    D3IL-avoiding benchmark, while the diffusion model loses it below its training budget
    (\autoref{sec:res:fewstep}). On D3IL-aligning analytic average-velocity matching improves with more
    evaluations and consistency-interpolated average-velocity matching does not.
  \item[RQ3 --- Constraints.] Endpoint projection is well-posed only with at least one guiding step
    (\autoref{sec:res:constraints:degenerate}). With several guiding steps it reduces time on D3IL-aligning
    and UAV-corridor and projection time on UAV-pillars; its constraint satisfaction is higher on D3IL-aligning,
    where it alone keeps every context free of violations, higher on UAV-pillars and lower on UAV-corridor (\autoref{sec:res:constraints}). D3IL-avoiding is
    already solved at one or two network evaluations, where endpoint projection has zero or at most one
    guiding step and no supported advantage, so per-step projection is the method there
    (\autoref{sec:res:avoiding:conclusion}).
  \item[RQ4 --- Transfer.] The lead of analytic average-velocity matching carries over from state to camera observations and from the
    manipulator to the quadrotor; the step from instantaneous-velocity matching to diffusion does not carry over to camera
    observations, where the two are level. On UAV-s-curve the tracking controller, not
    the plan, limits the result (\autoref{sec:res:uav:controller}).
\end{description}

\hole{Refine in v4.}

\section{Future Work}\label{sec:conc:future}

\hole{Written in v4. Outline material: \texttt{Writing/Working\_Space/future\_work/} (inferring the
constraint set from perception; changing constraints at test time).}
